\documentclass{article}
\usepackage[utf8]{inputenc}
\usepackage[english]{babel}
\usepackage{graphicx}
\usepackage{biblatex}
\usepackage{amssymb}
\usepackage{amsmath}
\usepackage{mathtools}
\usepackage{algorithm}
\usepackage[noend]{algpseudocode}
\usepackage{graphicx}
\usepackage{float}
\usepackage{scrextend}
\usepackage{hyperref}
\usepackage{caption}
\usepackage{subcaption}
 \usepackage{booktabs}
\usepackage[table,xcdraw]{xcolor}

\usepackage{amsthm}

\newtheorem{theorem}{Theorem}[section]
\newtheorem{corollary}{Corollary}[theorem]
\newtheorem{lemma}[theorem]{Lemma}

\theoremstyle{definition}
\newtheorem{definition}{Definition}[section]


\newcommand{\Enc}{\texttt{Enc}}
\newcommand{\Dec}{\texttt{Dec}}
\newcommand{\Gen}{\texttt{Gen}}

\newcommand{\M}{\mathcal{M}}
\renewcommand{\C}{\mathcal{C}}
\newcommand{\K}{\mathcal{K}}
\newcommand{\A}{\mathcal{A}}

\newcommand{\Prob}{\mathbb{P}}

\newcommand{\Oh}{\mathcal{O}}

\newcommand{\Int}{\mathbb{Z}}
\newcommand{\Nat}{\mathbb{N}}
\newcommand{\Reals}{\mathbb{R}}

\newcommand{\Expr}[2]{\texttt{Expr}^{\texttt{#1}}_{#2}}

\addtokomafont{labelinglabel}{\sffamily}

\usepackage[a4paper, total={6in, 8in}]{geometry}

\addbibresource{dissertation.bib}

\title{CS4796: Joint Senior Honours Project}
\author{\{gf38\} @ st-andrews.ac.uk}

\makeatletter
\def\BState{\State\hskip-\ALG@thistlm}
\makeatother

\begin{document}

\begin{titlepage}
    \begin{center}
        \vspace*{1cm}
        
        \textbf{Fully homomorphic encryption}
        
        \vspace{0.5cm}
        
        CS4796: Joint Senior Honours Project \\
                
        
        \vspace{1cm}
        
        \textbf{Gergely Flamich}

        \vspace{0.5cm}
        Supervisors:\\
        Dr Sophie Huczynska, Prof Steve Linton
        
        \vfill
        
        \includegraphics[width=0.5\textwidth]{img/uni_crest}
        
        \vspace{0.8cm}
        
        School of Computer Science\\
        University of St Andrews\\
        Scotland\\
        
    \end{center}
\end{titlepage}

\begin{abstract}
\end{abstract}
\paragraph{Declaration}
We declare that the material submitted for assessment is my own work except where credit is explicitly given to others by citation or acknowledgement. This work was performed during the current academic year except where otherwise stated. The main text of this project report is NN,NNN* words long, including project specification and plan. In submitting this project report to the University of St Andrews, I give permission for it to be made available for use in accordance with the regulations of the University Library. I also give permission for the report to be made available on the Web, for this work to be used in research within the University of St Andrews, and for any software to be released on an open source basis. I retain the copyright in this work, and ownership of any resulting intellectual property.

\newpage

\tableofcontents

\newpage


\section{Introduction}
\section{Mathematical Preliminaries}
\paragraph{}
Throughout this work some basic concepts will be used, the most important of
which will be given below.
\begin{definition}{Strings over an alphabet:}
  Let $A$ be a finite set of symbols, called an \textbf{alphabet} the strings
  over $A$ denoted $A^*$ is the set of all finite length sequences of
  symbols from $A$, including $\epsilon$, the empty string. Formally:
  \[
    A^* = \bigcup_{n\in \Nat \cup \{0\}} A^n.
  \]
  Here $\epsilon = () \in A^0$, the empty tuple. When talking about strings from $A^*$, instead of $s = (a_1, a_2, \hdots, a_n)
  \in A^*$, we will just write $a_1a_2\hdots a_n$.
\end{definition}
\paragraph{}
A natural property of strings is their length, i.e. the number of symbols they
are composed of. Thus, let $|\cdot| : A
\mapsto \Nat \cup \{0\}$ be defined as
\[
  \forall s \in A^* \quad |s| = n \Leftrightarrow s \in A^n
\]
\paragraph{}
A natural operation on strings is concatentaion. Formally, let $A$ be an
alphabet. Then, let $+ : A^* \times A^* \mapsto A^*$, such that
\[
  \forall s_1 = a_1\hdots a_n, s_2 = b_1 \hdots b_m \in A^*.\quad s_1 + s_2 =
  a_1\hdots a_nb_1 \hdots b_m \in A^*.
\]
\begin{definition}{Group:}
  Let $G$ be a non-empty set and $\circ: G \times G \mapsto G$ a binary operation on $G$.
  Then we will say that $(G, \circ)$ is a group if the following axioms hold:
  \begin{itemize}
  \item \textbf{G1:} $\forall a, b, c \in G$ we have $(a \circ b) \circ c = a
    \circ (b \circ c)$. (associativity)
  \item \textbf{G2:} $\exists 1 \in G. \,\forall a \in G$ we have $1 \circ a = a \circ 1 = a.$ (existence of an identity)
    \\
    It is easy to show that if such an element $1$ exists, then it
    is unique. We will call this unique element the \textbf{identity} of $G$.
    
  \item \textbf{G3:} $\forall a \in G\, \exists a^{-1} \in G$ such that $a \circ
    a^{-1} = a ^{-1} \circ a = 1$. (existence of inverses)
    \\
    It is easy to show that if this axiom holds
    then for every $a$ the corresponding $a^{-1}$ is unique. We will call this
    unique element the \textbf{inverse} of $a$. 
  \end{itemize}
\end{definition}
\paragraph{} If $G$ is finite, we call ($G, \circ$) a \textbf{finite group}.
\begin{definition}{Big-O notation:} Let $f, g: \Nat \mapsto \Reals $ be
  functions. Then we will write
  \[
    f(n) = \Oh(g(n)) \quad \text{(and say $f$ is of order big oh of $g$)}
  \]
  if and only if
  \[
    \exists M \in \Reals. \exists N \in \Nat.\quad \forall n \geq N \Rightarrow
    |f(n)| \leq M|g(n)|.
  \]
  Intuitively, this means that $f$ grows \textbf{at most} as fast as $g$.
\end{definition}
\section{A Very Brief History of Ciphers}
\paragraph{}
The history of ciphers goes back at least 2000 years. We will now examine the
cipher now named after Julius Ceasar, who used the following method to obfuscate
his correspondance for his adversaries:
\paragraph{Ceasar Cipher}: Take the message (written using symbols from the
Latin alphabet) that we wish to obfuscate. Replace each symbol with the one that
comes 3 places after it in the alphabet. For letters at the end where we could
not shift, we ``wrap around'' to the beginning of the alphabet and carry on
counting from there. For example
\[
  \textit{Alea iacta est} \quad\to\quad \textit{Dohd mdfzd hxz}.
\]
To decrypt a message, we simply shift backwards by 3 positions, e.g.
\[
  \textit{Zhqm, zmgm, zmgm} \quad\to\quad \textit{Veni, vidi, vici}.
\]
\paragraph{}
We can further generalise this concept to arrive at the definiton of
\textit{shift ciphers}, where instead of the fixing the shift to 3, we pick a
key $k \in \mathbb{N}$, and we then shift $k$ places forward (and backwards).
With shift ciphers in mind, we now formally defined some key concepts that will
be used throughout this work.
\begin{definition}{Cipher:}
  Let $\M, \C, \K_\Enc, \K_\Dec$ be alphabets. We will refer to $\M$
  as the \textbf{message space} and to $\C$ as the \textbf{cipher space}.
  A cipher $C$ over $\M$, $\K_\Enc$ and $\K_\Dec$ is a tuple $(\Gen, \Enc, \Dec)$, where
  \begin{itemize}
  \item $\Gen$ is the \textbf{key generation algorithm}, which outputs a key
    $(k_\Enc, k_\Dec) \in (\K_\Enc^* \times \K_\Dec^*)$, chosen according to some
    distribution. We refer to the set $\K\subseteq \K_\Enc^* \times \K_\Dec^*$ of all possible outputs of $\Gen$ as the
    \textbf{key space} of $C$. If for all outputs there is a polynomial time
    function $f$ such that $f(k_\Enc) = k_\Dec$, then we call $C$
    a \textbf{symmetric cipher} and instead of the tuple we just write $k$.
    Otherwise we call $C$ an \textbf{asymmetric cipher}.
    Since $\Gen$ is usually a probabilistic algorithm, we will denote generating
    its output by $k \leftarrow \Gen$ instead of $k = \Gen$ to emphasize the
    randomised nature of the function.
  \item $\Enc: \K_\Enc^*\times\M\mapsto \C$ is the \textbf{encryption algorithm}, which takes
    a encryption key $k$ and a message $m$ and outputs its encoding $c$. We will often
    refer to the message as the \textbf{plaintext} and to the encrypted message
    as \textbf{ciphertext}.
    \paragraph{Note:} $\Enc$ is not necessarily deterministic. When it is
    probabilistic, we will write $c \leftarrow \Enc_k(m)$ instead of $c = \Enc_k(m)$.
  \item $\Dec: \K_\Dec^*\times\C \mapsto \M$ is the \textbf{decryption algorithm}, which
    takes some decryption key $k$ and some ciphertext $c$ and outputs its
    correspnding plaintext $m$. 
    \paragraph{Note:} $\Dec$ is \textbf{always} deterministic (otherwise the
    scheme could never be assumed to be correct).
  \end{itemize}
  Finally, we require the correctness of $C$, concretely
  \[
    \forall m \in \M, \forall (k, k') \in \K. \quad \Dec(k', \Enc(k, m)) = m,
  \]
  i.e. that the cipher does not change its contents.
\end{definition}
\paragraph{Note:} For classical ciphers, we usually consider $\M = \C$ to be the
English (or Latin) alphabet. For modern crypto systems we will always work with
bitstrings, i.e. $\M$, $\C$ and $\K$ will be some subset of $\{0, 1\}^*$.
\paragraph{Equivalences of alphabets}
Consider any alphabet $\M = \{m_1, \hdots m_n\}$. It will be 
often convenient to consider the encoding of such an alphabet as something that
may be studied and manipulated more easily. Formally, by an encoding we mean a
bijection $f$ between our alphabet and some other set, where calculating values
of $f$ and $f^{-1}$ can both be done in polynomial time. Such an encoding may
be using ASCII for characters of the English alphabet or associating some
symbols with the elements of $\Int_n$. If there is such a bijection between
two alphabets, since one can be efficiently transformed into the other, we will
consider them equivalent.

\paragraph{Example}
Let us now revisit shift ciphers. In this case, for a set of $n$ symbols, it
will be easier to consider the encoding as elements of the additive group
$\Int_n$. First, it is easy to see that we are dealing with a
symmetric cipher, and since a shift by $k$ and any $k + xn, x \in \mathbb{N}$ is
going to be the same, $\K = \Int_n$. Then,
\begin{itemize}
\item $\Gen$: $\Gen$ is uniformly picking an element from $\Int_n$.
\item $\Enc$: We note that if we relabel our symbols to their
  indices (i.e. $m_i \to i$), then we can express our encryption function as
  \[
    \Enc(k, i) = i + k \mod n.
  \]
\item $\Dec$: Similarly, performing the above relabeling, we get
  \[
    \Dec(k, i) = i - k \mod n.
  \]
\end{itemize}
Now, checking the correctness of $C_{shift}$: Let $k, m \in \Int_n$
\[
  \Dec(k, \Enc(k, m)) = m + k \mod n - k \mod n = m + k - k \mod n = m \mod n = m.
\]
Since $k$ and $m$ were arbitrary, it holds $\forall m, k \in \Int_n$.
\paragraph{Kerckhoffs' principle}
Throughout the history of cryptography it has become evident that one must not
rely on the assumption that the adversary cannot obtain every detail about one's
crypto-system. This lead to the following important guiding principle in the
design of new schemes:
\begin{quotation}
The security of a scheme must lay in the key and not the scheme.
\end{quotation}
\section{Perfect Security}
In order for us to be able to analyse our schemes, we must have two well defined
notions: what does it mean that a scheme is \textit{secure} and how \textit{powerful} is
the adversary who wishes to break our scheme.
\paragraph{The adversary} We will start from a rather
paranoid, but very useful perspective: we will assume that the adversary has
unbounded computational power. This means that we will have to fend off attacks
that may involve computing arbitrary \textit{decidable} functions, no matter how
long they take to calculate. Also, by Kerckhoff's principle we assume that the
adversary has perfect knowledge of both our encryption and decryption functions.
\paragraph{Security} We will also impose a very stringent constraint on what we
consider secure.
\begin{definition}{Perfect Security}
  \label{def:perfect_security}
We say that a scheme $\Pi$ is perfectly secure if
\[
  \forall m \in \M, c, c' \in \C.\quad \Prob(M=m \,|\, C=c) = \Prob(M=m \,|\, C=c').
\]
\end{definition}
\paragraph{Indistinguishability} We will require that given any ciphertext and a uniformly
random string from the cipher alphabet, the adversary cannot
\textit{distinguish} the two strings whatsoever. To formalise this, we will
first need to define a \textbf{security experiment}:
\paragraph{}
Let $\Pi = (\Gen, \Enc, \Dec)$ be a scheme, and $\A$ be any adversary.
Then, we will define $\texttt{Expr}^{\texttt{eav}}_{\Pi, \A}$ as
follows:
\begin{enumerate}
\item $\A$ generates two messages $m_0$ and $m_1$ from the message space $\M$.
\item A random bit $b$ is chosen from $\{0, 1\}$.
\item Put $k \leftarrow \Gen()$ and $c = \Enc_k(m_b)$.
\item $\A$ is given $c$, and outputs $b'= \{0, 1\}$
\item The experiment outputs 1 if $b = b'$ and write $\Expr{eav}{\A, \Pi} = 1$,
  and 0 otherwise. If the output is 1, we say that $\A$ \textbf{succeeds}.
\end{enumerate}
Now, we are ready for our first definition of security:
\begin{definition}{Perfect Indistinguishability}
  \label{def:perfect_indistinguishability}
We say that a scheme $\Pi$ is perfectly indistinguishable iff for all adversaries $\A$
\[
  \Prob(\Expr{eav}{\A, \Pi} = 1) = \frac12.
\]
\end{definition}
\begin{theorem}
  \label{thm:perfect_equivalence}
A scheme $\Pi$ is perfectly secure if and only if it is perfectly indistinguishable.
\end{theorem}
\begin{proof}
  haha
\end{proof}
\begin{theorem}
A necessary condition for a scheme $\Pi = (\Gen, \Enc, \Dec)$ to be perfectly secure, is that $|\K|
\geq |\M|$.
\end{theorem}
\begin{proof}
  Assume $|\K| < |\M|$. We will show that this implies that perfect security is broken.\\
  The idea is to use the fact that if we are able to brute-force search through
  all the keys, we will learn some information about the keys and the messages.
  \paragraph{}
  For every $c \in \C$, let $\M_c = \{m \,|\, m = \Dec_k(c) k \in \K\}$.
  Note, that since $\Dec$ is a function, $|\M_c| \leq |K|$, as $\Dec$ may map two ciphertexts
  to the same plaintext for two different keys, but it can never map a
  ciphertext to two different plaintexts for the same key. Hence, $\M_c \subset
  \M$ and in particular $\M \setminus \M_c \neq \emptyset$.\\
  Now, fix $c \in \C$, and pick $m \in \M \setminus \M_c$. Since by definition
  $c$ cannot map to $m$, we have $\Prob(M=m | C = c) = 0$. But then
  \[
    \Prob(M=m) \neq 0 = \Prob(M=m | C=c)
  \]
  which violates perfect security, so our initial assumption must be false.
\end{proof}
\section{Computational Security}
\paragraph{}
Information theoretical security definions, although guaranteeing defense
against attacks of arbitrary strength, are problematic, as sharing and guarding
the keys are very cumbersome and in some cases impossible. In a sense, they are
also more powerful than what we need, since in the real world we do not have to
deal with attackers of unlimited computational power. Therefore, it would be
desirable to relax our notions of security that are more ``realistic'', while
the relaxation allows us to design much more efficient schemes.
In this section, we will discuss these relaxations, and 3 notions of security,
each being stronger than the previous. Since these definitions incorporate
assumptions about computational power and are no longer information theoretical,
we will call this approach \textbf{computational security}.
\section{Public-Key Crypto Schemes}
\subsection{An Example: RSA}
In this section we finally illustrate how the previously discussed ideas can be
combined into a concrete scheme. However, before we can do that we must make a
brief detour to prove some results we will be using later.
\subsubsection{Mathematical set-up}
\begin{theorem}{Order}
  Let $G$ be a finite group and $g \in G$. Then, there exist positive integers
  $n$ such that $g^n = 1$. We shall call the smallest such integer the order of
  $g$, denoted $|g|$ called the \textbf{order} of $g$.
\end{theorem}
\begin{proof}
  Since $G$ is finite, there must exist $n, m \in \Nat, n \neq m$ such that
  \[
    g^n = g^m.
  \]
  Without loss of generality, we may assume $n > m$. Then, multiplying both
  sides with the inverse of $g^m$, we get
  \[
    g^{n - m } = 1.
  \]
\end{proof}
\begin{theorem}
  \label{thm:elemorderdiv}
  Let $G$ be a finite group and $g \in G$. Then, the order of $g$ divides the
  order of $G$.
\end{theorem}
\begin{proof}
  Assume that $|g| \nmid |G|$. Then, $|G| = q|g| + r$ for $q$ a non-negative
  integer and $r$ a positive integer.
\end{proof}
\begin{corollary}
  Let $G$ be a finite group. Then,
  \[
    \forall g \in G\quad g^{|G|} = 1.
  \]
\end{corollary}
\begin{proof}
  Let $g \in G$. By Theorem \ref{thm:elemorderdiv} $|G| = q|g|$ for some $q \in
  \Nat$. Therefore,
  \[
    g^{|G|} = g^{q|g|} = (g^{|g|})^q = 1^q = 1.
  \]
\end{proof}
\begin{corollary}
  \label{thm:groupordpow}
  Let $G$ be a finite group. Then,
  \[
    \forall g \in G, p \in \Nat\quad g^p = g^{p \mod |G|}.
  \]
\end{corollary}
\begin{proof}
  \label{thm:grouppowmod}
  Let $g \in G, p \in \Nat$. Then, $p = q|G| + r$ for $q$ non-negative integer
  and $r$ a positive integer. Then, by the previous corollary $g^{|G|} = 1$, so
  \[
    g^p = g^{q|G| + r} = (g^{|G|})^qg^r = 1^qg^r= g^r = g^{p \mod |G|}.
  \]
\end{proof}
Before we may move on to the final corollary, we quickly make the following
observation:
\begin{theorem}
  Let $n \in \Nat$. The elements $x$ of $\Int_n \setminus \{0\}$ such that $\gcd(x, n) = 1$ form
  a group under multiplication with identity $1$ and inverse $a \in \Int_n$,
  where $ax + bn = 1$ from B\`ezout's identity. We denote this group by $\Int^*_n$
\end{theorem}
\begin{proof}
  It is very easy to check that the above indeed fulfils all group axioms.
\end{proof}
The order of $\Int_n$ is very important number-theoretically, so much so that it
is usually dentoted $\phi(n) = |\Int_n|$ and is called the \textbf{Euler phi-function}.
Now, we are ready to state the theorem of our interest:
\begin{theorem}{(Euler's Theorem)}\\
  Let $a, n \in \Nat$ such that $a < n$ and $\gcd(a, n) = 1$. Then,
  \[
    a^{\phi(n)} = 1 \mod n.
  \]
\end{theorem}
\begin{proof}
  By the above conditions $a \in \Int^*_n$. Then, the statement is a special
  case of Corollary \ref{thm:groupordpow}, where $G = \Int^*_n$.
\end{proof}
\subsubsection{The construction}

\section{Homomorphic Crypto Schemes}
\paragraph{}
At this point, we start to finally focus on the topic of the dissertation:
homomorphic schemes. Intuitively, these encapsulate the idea of encrypted
comptutation over encrypted data, in sense that the adversary may at most
observe the operation take place, but both the inputs and the output of it
remain secret. In this section we give the formal definition, construct an
example scheme known as the Paillier scheme. Finally, we close with a discussion
on the limitation of homomorphic schemes.
\subsection{Definition}
\subsection{Paillier's Scheme}
Paillier encryption is an additively homomorphic crypto scheme, that operates
over $\Int_{N^2}$ where $N = pq$ for some $p, q$ primes. In particular, it takes advantage of the fact that $\Int_{N^2}
\simeq \Int_{N} \times \Int_{N}^*$. 
\subsection{Limitations}
\section{Lattices}
\paragraph{}
In the previous section we have seen the motivating idea behind Gentry's scheme:
encrypted computation. In this section, we aim to provide a brief introduction
to the tool that helped the scheme come to life: lattices. We will define what a
lattice is and see some key definitions and theorems. We will close the section
with a construction of a simple lattice-based crypto scheme to illustrate how
the theory may be used to construct secure schemes with the appropriate choice
of a trapdoor.
\subsection{GGH}
\section{Fully Homomorphic Encryption}
\subsection{Dr Craig Gentry's Scheme}

\printbibliography

\end{document}